\section{Introduction}
\label{sec:intro}

Streaming 3D perception systems --- detectors and semantic mappers that run on
a live sensor feed~\cite{streamingperception,streaming3d,conceptfusion,conceptgraphs}
--- are consumed downstream as \emph{persistent objects}: a map, a scene
graph, a planner's object memory. What those consumers need is that one
physical object stay one object with one label. What streaming systems
actually produce flickers, because the prevailing emission rule is a per-frame
confidence threshold: every detection above a score cut-off enters the stream
immediately, independently of what was emitted a frame earlier. A re-detected
object arrives as a new identity, and an unstable open-vocabulary label changes
class along the way. Per-frame detection metrics do not see this, and so do not
penalise it. Temporal remedies exist --- confirmation counts, and confidence
accumulated along a track~\cite{cbmot,simpletrack,polymot,fastpoly} --- but
each also reduces how much a system emits, and detection-oriented metrics are
strongly coupled to output volume, so a comparison that does not hold that
volume fixed reports what a rule \emph{selects} and how much it \emph{removes}
together. The property is thus both unmeasured by the dominant metrics and
unattributed by the dominant comparison.

We study a \emph{retrospective confirmation} module in that setting, requiring
no retraining and no change to the perception model. Detections are linked into
tracks by a class-agnostic associator; a track becomes \emph{confirmed} only
after $N$ observations, and at that moment the module emits not only the
current detection but the track's withheld prefix, so a confirmed object's
trajectory is complete from its first observation. Nothing about a box's
coordinates, score, or label is modified, and the module runs at a single
operating point ($N{=}3$) with no per-dataset or per-metric tuning.

The attribution problem applies to our module as much as to any other, so we
evaluate under \emph{matched controls} that hold the output budget fixed while
the selection rule varies. A \emph{detection-budget-matched} control thresholds
the same frozen detector output to emit exactly as many detections as the
module; an \emph{identity-budget-matched} control keeps, per sequence, the
highest-scoring $K$ tracks, where $K$ is the number of identities the module
emits. To compare persistence against the other established temporal rule
rather than against score alone, a third control accumulates confidence along a
track in the manner of CBMOT~\cite{cbmot} and is held to both budgets. All
controls share the detector, the association step, and the evaluation code with
the module, differing only in the rule that selects what is emitted.

Under these controls, on the nuScenes validation set (150 scenes, a frozen
CenterPoint~\cite{centerpoint} detector, both sensor-frame and ego-motion-%
compensated world-frame association, confirmed tracks emitted
retrospectively), the module improves association accuracy (HOTA's AssA
component~\cite{hota}) over both matched controls in both association frames,
with bootstrap confidence intervals excluding zero, while class-aware tracking
accuracy AMOTA~\cite{nuscenes_tracking} --- a point estimate, since it does not
admit the same per-scene resampling --- rises consistently in sign and
magnitude. At the same output budget it \emph{reduces} per-frame detection
quality (Sec.~\ref{sec:detmatch}). The gain is bought rather than added.

Against confidence accumulation this becomes the paper's central observation:
at a matched budget the two temporal rules do not order one another but sit at
different points of the same trade-off. Persistence yields higher association
accuracy, accumulation retains the advantage on the detection- and
identity-oriented measures, and each of these differences has an interval
excluding zero. The separation is the same size whether emitted
\emph{detections} or emitted \emph{identities} are held fixed, which rules out
the simplest explanation --- that persistence gains association accuracy merely
by carrying fewer identities to confuse.

Because the module touches only emission timing, it transfers to any streaming
pipeline that re-observes objects. We apply the identical module, unchanged
and at the same $N$, to an indoor open-vocabulary instance-segmentation
pipeline that shares no detector, modality, or codebase with the outdoor one,
and evaluate it under the same control design. On ScanNet200 (312 validation
scenes, 3D instance-mask proposals labeled per frame by a 2D open-vocabulary
detector), the module reduces the number of frames in which an instance's
class label changes by 26.9\% relative to an identity-budget-matched control
of equal size, while leaving instance-segmentation AP essentially unchanged. The
temporal system has fewer label switches in 309 of 312 scenes and the paired
confidence interval excludes zero by a wide margin, whereas selecting the same
number of instances at random reproduces none of the reduction --- so the gain
is specific to temporal selection rather than to emitting fewer identities.

\noindent\textbf{Contributions.}
\begin{itemize}\itemsep2pt
\item \textbf{A training-free retrospective confirmation module.} An operator
combining the standard confirmation test with retrospective emission of the
confirmed prefix, applied to frozen detector output at a single operating
point, with no retraining, no architectural change, and no modification of box
coordinates, scores, or labels.
\item \textbf{The retrospective emission policy, isolated.} We show the
retrospective step is what completes the module: emitting causally instead
narrows \emph{where} the gain holds --- it survives in the sensor frame but
reverses in the world frame --- and retrospection's price is $N{-}1$ frames
of latency ($1.0$\,s at $2$\,Hz for $N{=}3$).
\item \textbf{A matched-budget comparison of two temporal selection rules.}
Confirmation and confidence accumulation are compared at an identical emission
budget and again at an identical identity budget, with paired per-sequence
statistics. Neither dominates: persistence buys association accuracy,
accumulation keeps the detection- and identity-oriented measures, and the
separation is unchanged under either matching.
\item \textbf{Controls that separate a rule from its suppression.} Both budgets
are also matched against a non-temporal score threshold, and indoors against
random selection, so what the comparisons credit is the selection rule rather
than the smaller volume it emits. The design is not specific to confirmation
and applies to any component that alters output volume.
\item \textbf{Cross-domain validation under that design.} Outdoors, the
module improves association accuracy and class-aware tracking accuracy under
both controls in both association frames, at a detection cost we quantify in
Sec.~\ref{sec:detmatch} rather than leave implicit. Indoors, the identical module
reduces class-label switching on an unrelated open-vocabulary
instance-segmentation pipeline, with segmentation quality unchanged.
\end{itemize}

\section{Related Work}
\label{sec:related}

\noindent\textbf{Streaming perception.} Streaming perception was formalised
by~\cite{streamingperception}, which scores predictions against the world at
the moment they are \emph{emitted}, folding latency and accuracy into streaming
AP (sAP); later work optimises detectors for this
regime~\cite{streamyolo}, and the 3D analogue processes LiDAR sectors as they
arrive rather than waiting for a full sweep~\cite{streaming3d}. This line
establishes our premise but corrects \emph{when} a prediction is scored: sAP
remains a per-frame average precision, so two systems emitting identical boxes
at identical latency but differing entirely in identity continuity and label
stability receive the same sAP. We hold latency and box geometry fixed and ask
what that family cannot see.

\noindent\textbf{Track lifecycle and confirmation.} Deciding when a candidate
track is real enough to report is among the oldest problems in multi-target
tracking. The classical radar literature settles it with \emph{track
confirmation} --- under $M$-out-of-$N$ history logic a track is confirmed once
it has been assigned a detection in $M$ of the last $N$ updates, under score
logic once its likelihood ratio crosses a sequential threshold --- and designs
confirmation, deletion, and track management together~\cite{blackmanpopoli}.
The same literature established that reporting may lag estimation:
\emph{retrodiction} revises earlier estimates once later measurements arrive,
in exchange for a bounded delay~\cite{retrodiction}. Tracking-by-detection
inherited the count-based half: SORT~\cite{sort} and DeepSORT~\cite{deepsort}
withhold a tentative track until matched in $N$ consecutive frames, and
AB3DMOT~\cite{ab3dmot} carries the same minimum-hits rule onto 3D boxes, while
CenterPoint~\cite{centerpoint} keeps a max-age deletion rule but reports tracks
from their first observation. None of this is ours: the confirmation test we
use is the textbook one at its conventional operating point. Our retrospective
emission likewise buys completeness with delay, as fixed-lag smoothing does,
and differs in what it revises --- it releases a confirmed track's withheld
prefix with coordinates, scores, and labels exactly as the detector produced
them, so the backward step changes set \emph{membership}, not state.

\noindent\textbf{Confidence in the track lifecycle.} Recent 3D tracking has
largely moved this decision from counting to confidence.
CBMOT~\cite{cbmot} replaces the minimum-hits criterion outright, propagating a
tracklet score by constant decay, fusing it probabilistically with each matched
detection score, and gating both output and termination on the refined score;
it reports consistent AMOTA gains over count-based lifecycles.
SimpleTrack~\cite{simpletrack} decomposes tracking-by-detection into
pre-processing, association, motion model, and lifecycle management, and names
\emph{output} as a lifecycle decision separate from birth and death, letting
low-score detections sustain a track through association without being emitted.
Poly-MOT~\cite{polymot} and Fast-Poly~\cite{fastpoly} pair category-specific
association with mixed count- and confidence-based termination and explicit
control over which trajectories are written out; Fast-Poly adopts CBMOT's
score-fusion rule directly. Confidence-based temporal management is therefore
established and well tuned, and we treat it as such: our contribution is not a
new confirmation or lifecycle mechanism, and our associator is deliberately
lightweight, so our absolute tracking scores are not comparable to these
systems'. What we take from this line is a comparison rather than a competitor.
Persistence-based confirmation and confidence accumulation are two temporal
rules that both decide which detections reach the output, and both reduce how
much is emitted; since detection-oriented metrics are strongly coupled to
output volume, a comparison that does not hold that volume fixed reports the
suppression and the rule together. Budget-matched comparison is routine when
scoring detectors at a fixed proposal count; we apply it to temporal selection,
holding fixed each of the two currencies such a rule spends --- emitted
detections and emitted identities --- and read what separates the two rules
under each.

\noindent\textbf{Open-vocabulary 3D perception.} OpenScene~\cite{openscene}
distils 2D vision--language features into 3D points and LERF~\cite{lerf} embeds
language into a radiance field; OpenMask3D~\cite{openmask3d} moves to instances
via per-mask CLIP features over class-agnostic 3D proposals, and
Open-YOLO~3D~\cite{openyolo3d} replaces that extraction with 2D
open-vocabulary boxes projected onto 3D masks. All four are \emph{offline},
consuming a completed reconstruction. ConceptFusion~\cite{conceptfusion} and
ConceptGraphs~\cite{conceptgraphs} make the setting incremental, merging
per-frame detections into persistent object nodes. A more recent line segments
3D instances fully online, computing proposals per frame:
EmbodiedSAM~\cite{embodiedsam} lifts per-frame 2D masks into 3D queries matched
across views by a learned similarity, OnlineAnySeg~\cite{onlineanyseg} merges
them into persistent instances zero-shot through voxel hashing,
and~\cite{autoseg3d} recasts online 3D segmentation as instance tracking with
long-term query association. Cross-frame identity is thus an acknowledged
design goal here, and in two of the three a trained objective, so we neither
call it neglected nor rest any distinction on being training-free --- one of
these systems is itself zero-shot. The gap is in what is measured: all are
scored by per-scene instance-segmentation AP, the tracking-framed system
included, so how often a persistent instance's \emph{class label} changes as
observations accumulate is never quantified. Our indoor experiment measures
that property, on a pipeline of this family, under a matched control.

\noindent\textbf{Evaluation metrics.} CLEAR MOTA~\cite{clearmot} counts FP, FN
and identity switches but is dominated by detection errors; IDF1~\cite{idf1}
rebalances toward identity by matching whole trajectories one-to-one.
HOTA~\cite{hota} was introduced because these conflate two capabilities, and
factorises into a detection accuracy DetA and an association accuracy AssA
whose geometric mean it is. That factorisation is what makes our result
legible: detection and association quality can be read separately over an
\emph{identical} box set, precisely the distinction a temporal selection rule
acts on. In 3D, nuScenes AMOTA~\cite{nuscenes_tracking} --- following AB3DMOT's
integral formulation~\cite{ab3dmot} --- averages MOTA over recall thresholds
and is class-aware, so it responds to labels the class-agnostic family ignores;
we report both. Video consistency metrics (STQ~\cite{step}, VPQ~\cite{vpq},
TETA~\cite{teta}) and open-vocabulary tracking benchmarks (TAO~\cite{tao},
OVTrack~\cite{ovtrack}) extend the family but require dense ground-truth tube
identities over a fixed taxonomy. We use these metrics where computable rather
than proposing a new one.
